\documentclass[10pt]{article}
\usepackage[letterpaper,left=0.45in,right=0.45in,top=0.45in,bottom=0.45in,headheight=14pt,headsep=8pt]{geometry}
\usepackage{multicol}
\usepackage{listings}
\usepackage{xcolor}
\usepackage{amsmath}
\usepackage{fancyhdr}

\ifdefined\XeTeXrevision
  \usepackage{fontspec}
  \setmonofont{Space Mono}
\else
  \usepackage{inconsolata}
\fi

\definecolor{codegray}{gray}{0.4}
\definecolor{docblue}{RGB}{15,55,120}
\definecolor{keywordblue}{RGB}{0,50,180}
\definecolor{boxbg}{gray}{0.96}

\lstset{
  language=C++,
  basicstyle=\fontsize{7.4}{8.4}\selectfont\ttfamily,
  commentstyle=\color{codegray}\itshape,
  keywordstyle=\color{keywordblue}\bfseries,
  stringstyle=\color{red!70!black},
  tabsize=2,
  showstringspaces=false,
  breaklines=true,
  frame=single,
  rulecolor=\color{gray!30},
  backgroundcolor=\color{boxbg},
  numbers=none,
  aboveskip=3pt,
  belowskip=3pt,
  columns=flexible,
  keepspaces=true,
  mathescape=true
}

\pagestyle{fancy}
\fancyhf{}
\fancyhead[L]{\textbf{\footnotesize ICPC TRD --- Guía}}
\fancyhead[R]{\footnotesize Pág. \thepage}
\renewcommand{\headrulewidth}{0.4pt}

\setlength{\columnsep}{16pt}
\setlength{\parindent}{0pt}
\setlength{\parskip}{3pt}
\setlength{\emergencystretch}{3em}

\begin{document}

\begin{center}
  {\Large\textbf{MANUAL TRD}}\\[2pt]
  {\small\textbf{Operación, Inventario de Plantillas y Reconocimiento de Problemas}}
\end{center}
\vspace{-4pt}
\hrule
\vspace{6pt}

\noindent\small Esta guía acompaña a \textit{plantilla.pdf}. Los algoritmos aparecen en el \textbf{mismo orden} que en la plantilla; el marcador \textbf{(0)/(1)} indica la base de indexación, igual que en los comentarios \textit{Uso[0]}/\textit{Uso[1]} de cada header. Cuando no aplica, no se marca.

\begin{multicols}{2}

\section*{0. Operación en Competencia}

\subsection*{0.1 Protocolo de la Máquina Única}
Hay 1 computadora para 3 personas; el recurso escaso es el tiempo de máquina.
\begin{itemize}\setlength{\itemsep}{1pt}
  \item \textbf{Papel primero:} nadie programa sin tener escritas las transiciones del DP, la condición de parada o la red de flujo con sus capacidades.
  \item \textbf{Regla de los 15 minutos:} si llevas 15 min atascado en un bug, cede la máquina y depura en papel mientras otro avanza.
  \item \textbf{Implementación por pares:} uno escribe y otro revisa el código en paralelo. Reduce errores de índices, typos y olvidos de casos borde.
  \item \textbf{Antes de cada submit:} corre el checklist de la sección 0.4 y casos de prueba extra (no solo los de la muestra).
\end{itemize}

\subsection*{0.2 Lectura y Gestión de Tiempo}
\begin{itemize}\setlength{\itemsep}{1pt}
  \item \textbf{Leer TODOS los problemas} en la primera pasada, antes de escribir código.
  \item No priorizar por ``los más resueltos'': el número de submits no mide dificultad. Un problema ignorado puede ser el más fácil.
  \item Anota por problema: idea, técnica, complejidad estimada y prioridad. Empieza por los fáciles confirmados.
  \item \textbf{Pausas cortas} (2--3 min) al estancarte para refrescar la mente; rota de problema y de máquina.
\end{itemize}

\subsection*{0.3 Roles}
\begin{itemize}\setlength{\itemsep}{1pt}
  \item \textbf{JuanJoe-18:} Sorting, Greedy, Geometría, Grafos, Range Queries y Strings. Identificación de estructuras y reducciones.
  \item \textbf{Botinix:} implementador principal; velocidad en templates y resolución de bugs difíciles.
  \item \textbf{JuanchoKing:} Chess, Matemáticas, Teoría de Números y Combinatoria. Reducciones abstractas.
\end{itemize}

\subsection*{0.4 Checklist Pre-Submit}
\begin{itemize}\setlength{\itemsep}{1pt}
  \item \textbf{Overflow:} si sumas o multiplicaciones pueden superar $2\times 10^9$, usa \textit{ll} y escribe \textit{1LL * a * b}.
  \item \textbf{INF:} para caminos mínimos usa \textit{LINF} (1e18); \textit{INF} (1e9) solo si los pesos lo permiten.
  \item \textbf{Múltiples testcases:} limpia estructuras globales solo hasta $N$; no reasignes con \textit{memset} de tamaños gigantes.
  \item \textbf{Indexación:} verifica la base (0)/(1) de cada estructura que uses (marcada en el inventario).
  \item \textbf{Casos borde:} $N=1,2$; grafo disconexo (¿imprime $-1$?); árbol en línea recta; cadena vacía o con un solo carácter; polígono degenerado; valores extremos (0, negativos).
  \item Imprime con \textit{'\textbackslash n'}, nunca \textit{endl}.
\end{itemize}

\subsection*{0.5 Stress Testing}
Cuando hay WA y un brute force confiable, verifica contra casos aleatorios pequeños. Bloque base:

\begin{lstlisting}[language=Python]
import random, subprocess

def gen():
    n = random.randint(1, 8)          # tamanos pequenos (brute rapido)
    a = [random.randint(1, 20) for _ in range(n)]
    return f"{n}\n{' '.join(map(str, a))}\n"

for it in range(2000):
    data = gen()
    sol   = subprocess.check_output("./sol",   input=data, text=True).strip()
    brute = subprocess.check_output("./brute", input=data, text=True).strip()
    if sol != brute:
        print("FALLO en:\n" + data)
        print("sol:", sol, "\nbrute:", brute)
        break
else:
    print("OK: no se encontraron diferencias")
\end{lstlisting}

\textbf{Uso:}
\begin{itemize}\setlength{\itemsep}{1pt}
  \item Compila tu solución y un brute force (misma lectura de entrada, misma salida) como \textit{sol} y \textit{brute} (en Windows, con \textit{sol.exe} y \textit{brute.exe}).
  \item \textbf{Qué variar en \textit{gen()}} según el tipo de problema:
  \begin{itemize}\setlength{\itemsep}{1pt}
    \item \textit{Arrays:} valores negativos y positivos, duplicados, extremos, tamaños 1 y 2.
    \item \textit{Grafos/árboles:} árboles aleatorios (nodo $i$ conectado a un padre $< i$), grafos con pesos, disconexos, con ciclos.
    \item \textit{Strings:} un solo carácter repetido, patrones que se solapan, cadenas cortas.
    \item \textit{Grillas:} $R$, $C$ pequeños, con obstáculos, bordes de tamaño 1.
    \item \textit{Flujo:} capacidades 0, cuellos de botella, demandas mínimas.
  \end{itemize}
  \item Si falla, reduce el caso a uno mínimo y depura con él.
  \item Para salidas de precisión (double), compara con tolerancia, no con igualdad.
  \item Si no aparece fallo en 2000 iteraciones, sube las iteraciones o cambia la semilla.
\end{itemize}

---

\section*{1. Main y Configuración (\textit{src/main.cpp})}
\begin{itemize}\setlength{\itemsep}{1pt}
  \item \textit{fastio}: \textit{sync\_with\_stdio} apagado.
  \item \textit{all}, \textit{rall}, \textit{pb}.
  \item typedefs \textit{ll}, \textit{ld}, \textit{pii}, \textit{vi}, \textit{vll}.
  \item constantes \textit{INF}, \textit{LINF}, \textit{MOD}, \textit{EPS}.
  \item \textit{rand\_val(a, b)}: entero aleatorio.
  \item pbds (\textit{ordered\_set}) incluido; \textit{solve()} con bucle de $T$ testcases.
\end{itemize}

\textbf{Tips y errores:}
\begin{itemize}\setlength{\itemsep}{1pt}
  \item Usa \textit{rand\_val} para bases aleatorias de hash y test aleatorio; el \textit{rng} local no colisiona con el de \textit{rq.hpp}.
  \item No uses \textit{endl} (flushea). Usa \textit{'\textbackslash n'}.
  \item Los pragmas de optimización están comentados; actívalos solo si el juez lo permite y el problema lo exige.
\end{itemize}

\section*{2. Grafos (\textit{graph.hpp})}

\textbf{Inventario (orden de la plantilla):}
\begin{itemize}\setlength{\itemsep}{1pt}
  \item[\textbf{Grillas}]
  \item \textit{dx4}, \textit{dy4}, \textit{dx8}, \textit{dy8} (1).
  \item \textit{isValid} (1).
  \item[\textbf{Internas}]
  \item \textit{Edge<T>}.
  \item[\textbf{Conjuntos}]
  \item \textit{DSU} (0): \textit{find}, \textit{unite}, \textit{component\_size}, \textit{members}, \textit{comps}, \textit{get\_components}.
  \item[\textbf{Estructura}]
  \item \textit{Graph<T>} (1): \textit{add\_directed\_edge}, \textit{add\_undirected\_edge}.
  \item[\textbf{Caminos mínimos}]
  \item \textit{Dijkstra<T>} (1): \textit{restore\_path}, \textit{invalid\_edges}.
  \item \textit{MultiStateDijkstra<T>} (1): multi-estado.
  \item \textit{ZeroOneBfs<T>} (1): pesos $\{0,1\}$.
  \item \textit{FloydWarshall<T>} (1).
  \item \textit{BellmanFord<T>} (1): ciclo negativo.
  \item \textit{Bipartite<T>} (1).
  \item[\textbf{Orden}]
  \item \textit{TopoSort<T>} (1): vacío si hay ciclo.
  \item \textit{DirectedCycle<T>} (1).
  \item \textit{UndirectedCycle<T>} (1).
  \item[\textbf{Componentes}]
  \item \textit{TarjanSCC<T>} (1).
  \item \textit{BridgeTree} (1): puentes, articulación, árbol de componentes 2-arista-conexas.
  \item[\textbf{MST}]
  \item \textit{Kruskal<T>} (1): devuelve índices de aristas.
  \item[\textbf{Caminos especiales}]
  \item \textit{EulerianCircuit<T>} (1).
  \item \textit{Hamiltonian} (1, $N \le 15$).
  \item \textit{FunctionalGraph} (1): k-ésimo sucesor, ciclos.
  \item \textit{DominatorTree} (1).
  \item[\textbf{LCA}]
  \item \textit{LCA<T>} (1): binary lifting.
  \item[\textbf{Booleano}]
  \item \textit{TwoSat} (1): cláusulas, implicación, xor, equivalencia.
  \item[\textbf{Emparejamiento}]
  \item \textit{BipartiteMatcher} (1): Kuhn.
  \item \textit{HopcroftKarp} (1).
  \item[\textbf{Clique}]
  \item \textit{MaxClique} (0, $N \le 60$).
  \item[\textbf{Flujo}]
  \item \textit{Dinic} (0): \textit{restore\_path}.
  \item \textit{DinicLB} (1): cotas inferiores.
  \item \textit{MCMF} (0): potenciales.
\end{itemize}

\textbf{Reconocimiento (patrón $\to$ algoritmo):}
\begin{itemize}\setlength{\itemsep}{1pt}
  \item ``Menor costo de $A$ a $B$'' $\to$ \textit{Dijkstra} (pesos $\ge 0$) o \textit{ZeroOneBfs} (solo $\{0,1\}$).
  \item ``Con hasta $K$ descuentos/boletos/combustible'' $\to$ \textit{MultiStateDijkstra}: estado $(nodo, recurso)$, respuesta = $\min_{rem} dist[t][rem]$.
  \item ``Costos negativos o detectar ciclo negativo'' $\to$ \textit{BellmanFord}.
  \item ``Todos los pares, $N \le 400$'' $\to$ \textit{FloydWarshall}.
  \item ``Conectar todo al mínimo costo'' $\to$ \textit{Kruskal}.
  \item ``Orden de tareas con dependencias'' $\to$ \textit{TopoSort} (si hay ciclo devuelve vacío).
  \item ``Recorrer cada arista una vez'' $\to$ Euleriano; ``cada nodo una vez'' $\to$ \textit{Hamiltonian} (solo $N \le 15$).
  \item ``Cada nodo apunta a un único nodo'' (funcional) $\to$ \textit{FunctionalGraph}: k-ésimo sucesor, ciclos, \textit{dist\_to\_cycle}.
  \item ``Asignar cada $X$ a un $Y$ sin repetir'' $\to$ \textit{BipartiteMatcher}/\textit{HopcroftKarp}. König: \textbf{min vertex cover = max matching}; \textbf{max independent set} $= N -$ matching. Dilworth: \textbf{máxima anticadena = min path cover} (con matching duplicando vértices).
  \item ``Proyectos con ganancia + costos + dependencias obligatorias'' (Project Selection) $\to$ min-cut: $S \to$ proyecto (ganancia), herramienta $\to T$ (costo), proyecto $\to$ herramienta ($\infty$); respuesta $= \Sigma$ganancias $-$ mincut.
  \item ``Variables binarias con restricciones por pares'' $\to$ \textit{TwoSat}. ``No ambas activas'' = $(A \Rightarrow \neg B) \wedge (B \Rightarrow \neg A)$; ``al menos una'' = $(\neg A \Rightarrow B) \wedge (\neg B \Rightarrow A)$; ``iguales'' = ambas implicaciones.
  \item ``Grafo implícito'' (estados como nodos, transiciones como aristas) $\to$ BFS / 0-1 BFS / Dijkstra.
  \item ``Minimizar el máximo / maximizar el mínimo / menor tiempo tal que...'' $\to$ \textbf{búsqueda binaria sobre la respuesta} con \textit{check} monótono (a menudo Dijkstra o greedy).
  \item \textbf{Grillas:} aplanar con \textit{id = r*C + c}; grilla 4-conexa es \textbf{bipartita} por $(r+c) \bmod 2$ (dominós $\to$ matching; König para bloquear celdas); capacidad por celda $\to$ dividir $u$ en $u_{in} \to u_{out}$ con capacidad $K$ y $u_{out} \to v_{in}$ con $\infty$.
\end{itemize}

\textbf{Errores y tips:}
\begin{itemize}\setlength{\itemsep}{1pt}
  \item \textit{Graph}, \textit{Dijkstra}, \textit{LCA}, \textit{TwoSat}, \textit{BridgeTree}, \textit{DinicLB}, \textit{DominatorTree} son \textbf{1-based}; \textit{DSU}, \textit{Dinic}, \textit{MCMF}, \textit{MaxClique} son \textbf{0-based}.
  \item \textit{DSU} es 0-based: si el problema da nodos $1..n$, usa \textit{DSU dsu(n+1)} (el nodo 0 queda aislado).
  \item \textit{MultiStateDijkstra}: \textit{rem} es el recurso \textbf{restante}; usar $K$ descuentos termina en \textit{rem=0}.
  \item \textit{invalid\_edges} requiere \textit{e.id != -1}; asigna ids al crear las aristas.
  \item \textit{Kruskal} devuelve índices de aristas, no de nodos.
  \item \textit{Dijkstra} con pesos negativos falla; usa \textit{BellmanFord}.
  \item \textit{MaxClique} es 0-based y exponencial: solo $N \le 60$ con bitsets.
  \item \textit{EulerianCircuit}: usa \textit{e.id} para no repetir aristas; en no dirigido el circuito tiene $E/2+1$ vértices.
\end{itemize}

\section*{3. Árboles (\textit{tree.hpp})}

\textbf{Inventario (orden de la plantilla):}
\begin{itemize}\setlength{\itemsep}{1pt}
  \item[\textbf{Internas}]
  \item \textit{TreeEdge<T>}.
  \item[\textbf{Árbol base}]
  \item \textit{Tree<T>} (1): raíz, euler \textit{tin}/\textit{tout}, \textit{sz}, \textit{depth}, \textit{get\_diameter} con camino, \textit{get\_centers}.
  \item \textit{all\_farthest\_distances}: distancia máxima por nodo.
  \item \textit{BinaryLifting<T>} (1): \textit{lca}, \textit{kth\_ancestor}, \textit{path\_sum}, \textit{path\_max}, \textit{path\_min}, \textit{dist\_weighted}, \textit{kth\_node\_on\_path}, \textit{is\_ancestor}.
  \item \textit{EulerTourTree<T>} (1): mapea subárbol a \textit{[tin, tout]}.
  \item \textit{HLD} (1): \textit{values\_on\_edges} nativo, \textit{process\_path}, \textit{process\_subtree}.
  \item[\textbf{Centroid}]
  \item \textit{CentroidPathProcessing} (1): Query-then-Add.
  \item \textit{CentroidTree<T>} (1): distancia al nodo activo más cercano.
  \item \textit{CentroidBinarySearch} (1): nodo oculto interactivo.
  \item[\textbf{Subárboles y caminos}]
  \item \textit{DSUOnTree} (1): small-to-large.
  \item \textit{VirtualTree} (1): con lambdas de LCA.
  \item \textit{TreeHashing} (1): isomorfismo.
  \item \textit{RerootingDP} (1).
  \item \textit{FastLCA} (1): Euler+RMQ $O(1)$.
  \item \textit{MoOnTrees} (1).
\end{itemize}

\textbf{Reconocimiento:}
\begin{itemize}\setlength{\itemsep}{1pt}
  \item ``Ancestro común / distancia en árbol'' $\to$ \textit{BinaryLifting} ($O(\log N)$) o \textit{FastLCA} ($O(1)$ si hay muchas consultas).
  \item ``Consultas sobre caminos (suma/mín/máx/asignación)'' $\to$ \textit{HLD} + \textit{LazySegTree}; con pesos en aristas usa \textit{values\_on\_edges = true}.
  \item ``Consultas sobre subárboles'' $\to$ \textit{EulerTourTree} + Fenwick/SegTree.
  \item ``Para cada nodo como raíz, calcula $X$'' $\to$ \textit{RerootingDP} (merge + extend + finalize).
  \item ``Contar/optimizar caminos por el árbol (ej. exactamente $K$ aristas)'' $\to$ \textit{CentroidPathProcessing}.
  \item ``Actualizar y consultar el nodo activo más cercano / suma de distancias'' $\to$ \textit{CentroidTree} (inyecta la función de distancia).
  \item ``Encontrar un nodo oculto con consultas interactivas'' $\to$ \textit{CentroidBinarySearch}.
  \item ``Estadísticas de colores por subárbol'' $\to$ \textit{DSUOnTree}.
  \item ``DP sobre un subconjunto de nodos (con LCA de por medio)'' $\to$ \textit{VirtualTree}.
  \item ``Frecuencias de colores/valores en caminos $u$--$v$'' $\to$ \textit{MoOnTrees}.
  \item ``¿Dos árboles son isomorfos?'' $\to$ \textit{TreeHashing}.
  \item Diámetro y centros $\to$ \textit{Tree.get\_diameter/get\_centers}.
\end{itemize}

\textbf{Errores y tips:}
\begin{itemize}\setlength{\itemsep}{1pt}
  \item Todo en \textit{tree.hpp} es \textbf{1-based}.
  \item \textit{Tree} requiere \textit{init(root)} antes de \textit{tin} y \textit{tout}; \textit{all\_farthest\_distances} usa el diámetro y asume el árbol inicializado.
  \item \textit{HLD} con \textit{values\_on\_edges}: el LCA se excluye automáticamente; no consultes el nodo raíz.
  \item \textit{RerootingDP}: \textit{merge} debe ser asociativo/conmutativo para el paso top-down (usa prefijos/sufijos si no es invertible).
  \item \textit{CentroidPathProcessing}: modifica \textit{PathState}, \textit{GlobalState}, \textit{extend}, \textit{query\_and\_merge}, \textit{add\_to\_global}, \textit{clear\_global}.
  \item \textit{CentroidTree}: modifica \textit{State}, \textit{NEUTRAL}, \textit{merge}, \textit{combine}.
  \item \textit{MoOnTrees}: implementa \textit{toggle(u)} antes de \textit{solve}.
  \item \textit{VirtualTree}: recibe \textit{get\_lca}, \textit{is\_anc}, \textit{get\_dist}, y orden por \textit{tin} como lambdas.
\end{itemize}

\section*{4. Range Queries (\textit{rq.hpp})}

\textbf{Inventario (orden de la plantilla):}
\begin{itemize}\setlength{\itemsep}{1pt}
  \item \textit{SparseTable} (0): consultas estáticas idempotentes (min/max/gcd).
  \item \textit{FenwickTree} (1): sumas + update puntual.
  \item \textit{IterativeSegTree} (0): suma/min/max + update puntual.
  \item \textit{RecursiveSegTree} (0): \textit{Node}+\textit{merge}, \textit{find\_first}.
  \item \textit{MaxSubarraySegTree} (0): máxima suma de subsegmento contiguo.
  \item \textit{LazySegTree} (0): \textit{add\_range}, \textit{set\_range}, \textit{get}.
  \item \textit{DynamicSegTree} (1..max): sparse sin compresión.
  \item \textit{PersistentSegTree} (0): versiones históricas.
  \item \textit{FenwickTree2D} (1): submatrices + update puntual.
  \item \textit{MergeSortTree} (0): conteo $\le k$ en rango + update.
  \item \textit{SqrtDecomposition} (0).
  \item \textit{Treap} (0): orden estadístico.
  \item \textit{ImplicitTreap} (0): secuencia con reverse.
  \item \textit{Mo} (0): consultas offline.
  \item \textit{MoWithUpdates} (0): Mo 3D.
  \item \textit{RollbackMo} (0): solo add.
\end{itemize}

\textbf{Reconocimiento:}
\begin{itemize}\setlength{\itemsep}{1pt}
  \item RMQ estático $\to$ \textit{SparseTable}.
  \item Suma con updates puntuales $\to$ \textit{Fenwick} (más rápido y simple que el segment tree).
  \item Operación no idempotente o con updates $\to$ segment tree; si la operación es compleja, redefine \textit{Node}+\textit{merge} en \textit{RecursiveSegTree}.
  \item ``Máxima suma de subsegmento en un rango'' $\to$ \textit{MaxSubarraySegTree} (nodo con \textit{sum, pref, suff, ans}).
  \item ``Suma/asignación sobre un rango + consulta'' $\to$ \textit{LazySegTree}.
  \item ``Índices enormes sin comprimir'' $\to$ \textit{DynamicSegTree}.
  \item ``Consultas sobre versiones anteriores'' (ej. k-ésimo menor en rangos) $\to$ \textit{PersistentSegTree}.
  \item ``Contar elementos $\le k$ en un rango (con updates)'' $\to$ \textit{MergeSortTree}.
  \item ``Consultas offline de rangos'' $\to$ \textit{Mo}; con updates $\to$ \textit{MoWithUpdates}; si solo se agrega (sin remover fácil) $\to$ \textit{RollbackMo}.
  \item ``Secuencia dinámica con insert/borrar/revertir'' $\to$ \textit{ImplicitTreap}.
  \item ``k-ésimo menor de un multiset dinámico'' $\to$ \textit{Treap}.
\end{itemize}

\textbf{Errores y tips:}
\begin{itemize}\setlength{\itemsep}{1pt}
  \item \textbf{Fenwick:} \textit{idx = 0} entra en \textbf{bucle infinito} (\textit{0 \& -0 = 0}). Si tus datos son 0-based, haz \textit{ft.add(x+1, v)} y \textit{ft.query(l+1, r+1)}.
  \item 0-based: \textit{SparseTable}, \textit{IterativeSegTree}, \textit{RecursiveSegTree}, \textit{LazySegTree}, \textit{MaxSubarraySegTree}, \textit{PersistentSegTree}.
  \item 0-based: \textit{MergeSortTree}, \textit{SqrtDecomposition}, \textit{Mo}, \textit{MoWithUpdates}, \textit{RollbackMo}, \textit{ImplicitTreap}. 1-based: \textit{FenwickTree}, \textit{FenwickTree2D}. \textit{DynamicSegTree} usa $1 \dots$ max.
  \item \textit{LazySegTree}: \textit{update(type 1=add, 2=set)}; tras un \textit{set\_range}, un \textit{add\_range} posterior suma sobre el valor seteado.
  \item \textit{PersistentSegTree}: \textit{update} devuelve la nueva raíz; guarda cada versión.
  \item \textit{Mo}: las queries se ordenan internamente; \textit{add}/\textit{remove} deben ser $O(1)$.
  \item \textit{Treap.kth\_element} es 0-based y devuelve \textit{LINF} fuera de rango.
\end{itemize}

\textbf{Variaciones:}
\begin{itemize}\setlength{\itemsep}{1pt}
  \item Desacoplar \textit{Node}+\textit{merge} (max subarray ya en plantilla): nunca toques los índices del árbol.
  \item Lazy para otra operación: edita \textit{merge} y cómo \textit{apply\_add}/\textit{apply\_set} combinan el valor del nodo.
  \item Segment tree + Euler (subárboles) o + \textit{HLD} (caminos): los rangos que produce el árbol alimentan la estructura 1D.
\end{itemize}

\section*{5. Strings (\textit{string.hpp})}

\textbf{Inventario (orden de la plantilla):}
\begin{itemize}\setlength{\itemsep}{1pt}
  \item[\textbf{Utilidades}]
  \item \textit{to\_lower}, \textit{to\_upper}, \textit{trim}, \textit{replace\_all}, \textit{split}, \textit{join}.
  \item[\textbf{Patrones y hashing}]
  \item \textit{prefix\_function} (KMP).
  \item \textit{z\_function}.
  \item \textit{remove\_all\_occurrences} (KMP con pila).
  \item \textit{rabin\_karp} (0): hash simple.
  \item \textit{manacher}: radios de palíndromos.
  \item \textit{StringHash} (0): doble módulo, bases configurables o aleatorias.
  \item[\textbf{Estructuras}]
  \item \textit{Trie}.
  \item \textit{AhoCorasick} (0): multi-patrón.
  \item \textit{SuffixAutomaton} (0).
  \item \textit{SuffixArray} (0): con LCP RMQ.
  \item \textit{PalindromicTree} (0).
\end{itemize}

\textbf{Reconocimiento (cuál usar):}
\begin{itemize}\setlength{\itemsep}{1pt}
  \item \textit{KMP (pi)}: periodo de una cadena, prefijo que es sufijo, autómata de un patrón.
  \item \textit{Z}: para cada sufijo, LCP con el inicio de la cadena; más intuitivo que KMP para modificar.
  \item \textit{remove\_all\_occurrences}: eliminar un patrón incluyendo ocurrencias encadenadas (ej. ``aaa'' con ``aa'').
  \item \textit{StringHash}: comparar substrings en $O(1)$; con \textit{true} usa bases aleatorias (inmune a colisiones forzadas en jueces adversariales).
  \item \textit{rabin\_karp}: posiciones de un patrón; base fija $\Rightarrow$ preferir \textit{StringHash} si hay riesgo de colisiones.
  \item \textit{AhoCorasick}: \textbf{múltiples patrones} en un texto (conteo y primera aparición).
  \item \textit{SuffixAutomaton}: substrings distintos, ocurrencias de un patrón, primera aparición.
  \item \textit{SuffixArray}: sufijos ordenados + LCP entre cualquier par en $O(1)$ (repeticiones, substrings comunes, conteos con lcp).
  \item \textit{Manacher}: radio del palíndromo más largo en cada centro.
  \item \textit{PalindromicTree} (eertree): todos los palíndromos distintos, sus ocurrencias y primera posición.
  \item \textit{Trie}: autómatas de prefijos, inserciones/búsquedas.
\end{itemize}

\textbf{Errores y tips:}
\begin{itemize}\setlength{\itemsep}{1pt}
  \item \textit{StringHash} y \textit{rabin\_karp} asumen minúsculas (\textit{s[i]-'a'+1}); adapta la constante si hay otros caracteres.
  \item \textit{manacher} devuelve radios; el máximo radio es la longitud del palíndromo más largo.
  \item \textit{SuffixArray} agrega internamente un \textit{'\$'} a la cadena; \textit{p[0]} es el sufijo ``\$''.
  \item \textit{PalindromicTree}: \textit{count\_palindromes} son los distintos; llama \textit{compute\_occurrences()} antes de leer \textit{occ}.
  \item \textit{StringHash(s, true)} genera bases distintas por proceso: no compares hashes de instancias diferentes.
\end{itemize}

\section*{6. Programación Dinámica (\textit{dp.hpp})}

\textbf{Inventario (orden de la plantilla):}
\begin{itemize}\setlength{\itemsep}{1pt}
  \item \textit{LiChaoTree}: mínimo de $m \cdot x + c$ (x arbitraria).
  \item \textit{dnc\_dp}: divide \& conquer para $dp[i][j] = \min_{k<j}(dp[i-1][k] + cost(k+1, j))$ con decisión monótona.
  \item \textit{knuth\_dp}: optimización de Knuth para DP de intervalos.
  \item \textit{DigitDP}.
  \item \textit{sos\_dp}: suma sobre submáscaras.
  \item \textit{lis}: con reconstrucción.
  \item \textit{bitset\_knapsack}.
  \item \textit{WQS}: Alien's trick.
  \item \textit{CHTDeque}: pendientes monótonas.
\end{itemize}

\textbf{Marco teórico (DP = caminos sobre un DAG):}
\begin{itemize}\setlength{\itemsep}{1pt}
  \item Estado = nodo; transición = arista; DP = caminos sobre el DAG inducido. Mínimo costo = camino más corto; máximo beneficio = camino más largo; conteo = número de caminos.
  \item Ante un DP ``duro'', responde en papel: (1) \textbf{eje de avance}: ¿qué variable avanza monótono (índice, suma, tiempo, bits)?; (2) \textbf{memoria mínima}: ¿qué información del pasado condiciona el futuro?; (3) \textbf{decisiones legales}: ¿qué transiciones e invariantes hay?
  \item \textbf{Permutaciones vs combinaciones:} contar composiciones que suman $S$ con 1D ($dp[s] = \sum_c dp[s-c]$, distingue orden) vs 2D ($dp[i][s]$ con índice de moneda, orden fijo, colapsa combinaciones).
\end{itemize}

\textbf{Reconocimiento (qué optimización):}
\begin{itemize}\setlength{\itemsep}{1pt}
  \item ``$dp[i] = \min_j (m_j \cdot x_i + b_j)$'' $\to$ \textit{LiChaoTree} (x arbitraria) o \textit{CHTDeque} (pendientes y $x$ monótonos).
  \item ``Particionar en $K$ grupos / $dp[i][j] = \min_k (dp[i-1][k] + C(k,j))$ con $opt$ monótono'' $\to$ \textit{dnc\_dp}.
  \item ``Intervalos: $dp[i][j] = \min_k(dp[i][k] + dp[k][j]) + cost$'' $\to$ \textit{knuth\_dp}.
  \item ``Contar números que cumplen condiciones de dígitos'' $\to$ \textit{DigitDP} (recuerda \textit{tight} y \textit{lz}).
  \item ``Suma sobre subconjuntos de máscaras'' $\to$ \textit{sos\_dp} ($O(N 2^N)$).
  \item ``Subsecuencia creciente más larga (con reconstrucción)'' $\to$ \textit{lis}.
  \item ``¿Subconjunto suma exactamente $S$?'' $\to$ \textit{bitset\_knapsack} (rápido, $S \le 10^5$ fijo).
  \item ``Exactamente $K$ elementos/particiones con función convexa'' $\to$ \textit{WQS} (Alien's trick).
  \item ``Para cada nodo como raíz'' $\to$ \textit{RerootingDP} (tree.hpp).
  \item \textbf{Greedy:} si no hay subestructura que exija DP (actividades, intervalos, monedas canónicas), ordena y elige localmente; verifica con un contraejemplo antes de descartar.
\end{itemize}

\textbf{Errores y tips:}
\begin{itemize}\setlength{\itemsep}{1pt}
  \item \textit{dnc\_dp}: \textit{cost} debe evaluarse en $O(1)$ (usa prefijos); para particionar en $K$ grupos, \textit{dp\_prev[0] = 0}; 0-based.
  \item \textit{knuth\_dp}: requiere quadrangle inequality en \textit{cost}.
  \item \textit{DigitDP}: \textit{memo} se resetea en cada \textit{solve}; \textit{lz} evita contar ceros a la izquierda.
  \item \textit{WQS}: \textit{check(mid).se >= K} decide la búsqueda.
  \item \textit{CHTDeque.query(x)} exige $x$ no decreciente; usa \textit{query\_bs} si $x$ es arbitraria.
  \item \textit{bitset\_knapsack}: \textit{MAX\_SUM} está fijo en $10^5$; si el objetivo es mayor, redimensiona la constante.
\end{itemize}

\section*{7. Matemáticas (\textit{math.hpp})}

\textbf{Inventario (orden de la plantilla):}
\begin{itemize}\setlength{\itemsep}{1pt}
  \item \textit{modpow}.
  \item \textit{ext\_gcd}.
  \item \textit{modinv}.
  \item \textit{crt}: dos congruencias.
  \item \textit{solve\_crt\_array}.
  \item \textit{Sieve}: primos, \textit{spf}, \textit{mu}, \textit{phi}, \textit{factorize}.
  \item \textit{Combinatorics}: \textit{nCr}, \textit{catalan}.
  \item \textit{gauss}: sistemas lineales (double).
  \item \textit{multiply} (FFT): polinomios enteros.
  \item \textit{multiply\_mod} (NTT).
  \item \textit{nim\_game}.
  \item \textit{sum\_n}, \textit{sum\_squares}, \textit{sum\_cubes}, \textit{geom\_sum}.
  \item \textit{Matrix}: multiplicación y potencia modular.
  \item \textit{berlekamp\_massey}.
  \item \textit{PollardRho}: \textit{is\_prime} (Miller-Rabin), \textit{factorize} hasta $2^{64}$.
  \item \textit{Lucas}.
  \item \textit{xor\_convolution}, \textit{or\_convolution}, \textit{and\_convolution} (FWHT).
  \item \textit{floor\_sum}.
\end{itemize}

\textbf{Reconocimiento:}
\begin{itemize}\setlength{\itemsep}{1pt}
  \item ``Recurrencia lineal grande / número de caminos en $k$ pasos'' $\to$ \textit{Matrix.power}.
  \item ``Conjeturar la recurrencia de una secuencia'' $\to$ \textit{berlekamp\_massey}.
  \item ``$nCr$ con $n$ enorme y módulo primo pequeño'' $\to$ \textit{Lucas}.
  \item ``Multiplicar polinomios'' $\to$ \textit{multiply} (FFT, enteros) o \textit{multiply\_mod} (NTT). ``Convoluciones XOR/OR/AND'' $\to$ FWHT.
  \item ``Sistema de congruencias'' $\to$ \textit{crt}/\textit{solve\_crt\_array}.
  \item ``Factorizar enteros de hasta $2^{64}$ / ¿es primo?'' $\to$ \textit{PollardRho}.
  \item ``Möbius / totient / factorizar muchos valores'' $\to$ \textit{Sieve}.
  \item ``Juego de Nim'' $\to$ xor de pilas (\textit{nim\_game}).
  \item ``Sumatorias cerradas'' $\to$ \textit{sum\_n} etc.; ``$\sum \lfloor (ai+b)/m \rfloor$'' $\to$ \textit{floor\_sum}.
  \item ``Sistema de ecuaciones lineales'' $\to$ \textit{gauss}.
\end{itemize}

\textbf{Errores y tips:}
\begin{itemize}\setlength{\itemsep}{1pt}
  \item \textit{modinv} exige $\gcd(a, m) = 1$ (\textit{assert}).
  \item \textit{crt} con congruencias inconsistentes devuelve $\{-1, -1\}$; verifícalo.
  \item NTT requiere tamaño potencia de 2 y que el módulo tenga la raíz (1e9+7 usa raíz primitiva 3).
  \item FFT usa \textit{double}: para coeficientes modulares muy grandes prefiere \textit{multiply\_mod}.
  \item \textit{Matrix} multiplica en $O(r \cdot c^2)$: cuidado con matrices grandes.
  \item \textit{Lucas} preprocesa factoriales hasta $p-1$ mód $p$.
\end{itemize}

\section*{8. Geometría (\textit{geometry.hpp})}

\textbf{Inventario (orden de la plantilla):}
\begin{itemize}\setlength{\itemsep}{1pt}
  \item \textit{cmp}: comparación con EPS.
  \item \textit{Point<T>}: operadores + orden lexicográfico.
  \item \textit{dot}, \textit{cross}, \textit{cross(p,a,b)}, \textit{norm2}, \textit{norm}.
  \item \textit{sgn}, \textit{orient}.
  \item \textit{on\_segment}.
  \item \textit{segment\_intersect}.
  \item \textit{polygon\_area\_2}: doble del área.
  \item \textit{boundary\_points}: Pick.
  \item \textit{point\_in\_polygon}.
  \item \textit{convex\_hull}: preserva colineales.
  \item \textit{closest\_pair}.
  \item \textit{Event1D}, \textit{Event2D}: sweep line.
  \item \textit{polar\_sort}.
  \item \textit{dist\_to\_line}.
  \item \textit{dist\_to\_segment}.
  \item \textit{line\_intersect}.
  \item \textit{rotating\_calipers}: diámetro.
  \item \textit{Circle}.
  \item \textit{circle\_line\_intersect}.
  \item \textit{circle\_circle\_intersect}.
  \item \textit{Halfplane}.
  \item \textit{halfplane\_intersection}.
  \item \textit{Point3D}.
  \item \textit{dist\_point\_plane}, \textit{dist\_point\_line}.
  \item \textit{line\_plane\_intersect}.
\end{itemize}

\textbf{Reconocimiento:}
\begin{itemize}\setlength{\itemsep}{1pt}
  \item ``Orientación / colinealidad'' $\to$ \textit{cross}, \textit{orient} (enteros).
  \item ``¿Punto dentro de un polígono? / área / puntos del retículo'' $\to$ \textit{point\_in\_polygon}, \textit{polygon\_area\_2}, \textit{boundary\_points} (Pick: $A = I + B/2 - 1$).
  \item ``Cáscara convexa / envolvente'' $\to$ \textit{convex\_hull}; ``diámetro de un convexo'' $\to$ \textit{rotating\_calipers}.
  \item ``Par de puntos más cercano'' $\to$ \textit{closest\_pair}.
  \item ``Intersección de segmentos/rectas/círculos'' $\to$ \textit{segment\_intersect}, \textit{line\_intersect}, \textit{circle\_*}.
  \item ``Región factible de restricciones lineales (semiplanos)'' $\to$ \textit{halfplane\_intersection}.
  \item ``Intervalos que se solapan / área de unión de rectángulos'' $\to$ sweep line con \textit{Event1D}/\textit{Event2D} + segment tree.
  \item ``Ordenar puntos por ángulo'' $\to$ \textit{polar\_sort}.
\end{itemize}

\textbf{Errores y tips:}
\begin{itemize}\setlength{\itemsep}{1pt}
  \item Usa enteros (\textit{ll}) mientras puedas; con doubles usa \textit{cmp} y \textit{EPS}.
  \item \textit{convex\_hull} preserva colineales en el borde (usa \textit{cross < 0}); para cáscara estricta cambia a \textit{cross <= 0}.
  \item \textit{polygon\_area\_2} devuelve \textbf{el doble} del área (divide entre 2).
  \item \textit{line\_intersect} no es válida si las rectas son paralelas.
  \item \textit{halfplane\_intersection} devuelve vacío si la región es ilimitada o no existe; el lado válido es la izquierda de la recta orientada.
  \item \textit{point\_in\_polygon} devuelve 0 (fuera), 1 (borde), 2 (dentro).
\end{itemize}

\section*{9. Bits (\textit{bits.hpp})}

\textbf{Inventario (orden de la plantilla):}
\begin{itemize}\setlength{\itemsep}{1pt}
  \item \textit{check\_bit}, \textit{set\_bit}, \textit{clear\_bit}, \textit{toggle\_bit}, \textit{lsb} (bit $i$ 0-based).
  \item \textit{count\_bits}, \textit{count\_bits\_ll}, \textit{highest\_bit}, \textit{ctz}.
  \item \textit{iterate\_submasks}: itera submáscaras (sin la vacía).
  \item \textit{gospers\_hack}: máscaras con $k$ bits.
  \item \textit{XorBasis}: base lineal xor.
  \item \textit{BitTrie}: máximo xor con borrado.
\end{itemize}

\textbf{Reconocimiento:}
\begin{itemize}\setlength{\itemsep}{1pt}
  \item ``Máximo xor de un subconjunto / ¿se puede formar $k$?'' $\to$ \textit{XorBasis}.
  \item ``Máximo xor de un valor contra un conjunto (dinámico)'' $\to$ \textit{BitTrie}.
  \item ``DP sobre submáscaras'' $\to$ \textit{iterate\_submasks} (total $O(3^N)$) o \textit{sos\_dp} ($O(N 2^N)$).
  \item ``Subconjuntos de tamaño exacto $k$'' $\to$ \textit{gospers\_hack}.
  \item ``Comparar multiconjuntos de rangos'' $\to$ \textit{ZobristHash} (extra).
\end{itemize}

\textbf{Errores y tips:}
\begin{itemize}\setlength{\itemsep}{1pt}
  \item Escribe \textit{1LL << i} (no \textit{1 << i}) con $i \ge 31$.
  \item \textit{BitTrie} usa 30 bits por defecto (hasta $\approx 10^9$); sube a 60 para valores hasta $10^{18}$.
  \item \textit{iterate\_submasks} no incluye la submáscara vacía (procesa \textit{sub = 0} aparte si la necesitas).
  \item \textit{highest\_bit} devuelve $-1$ para 0; \textit{ctz} devuelve 64 para 0.
\end{itemize}

\section*{10. Extra (\textit{extra.hpp})}

\textbf{Inventario (orden de la plantilla):}
\begin{itemize}\setlength{\itemsep}{1pt}
  \item \textit{ordered\_set} (0): pbds, k-ésimo menor, \textit{order\_of\_key}.
  \item \textit{ordered\_multiset} (0): pbds, permite duplicados.
  \item \textit{custom\_hash}: splitmix64 anti-colisión.
  \item \textit{MeetInTheMiddle}.
  \item \textit{ZobristHash} (0).
  \item \textit{BitsetReachability} (1).
  \item \textit{compress\_coords} (0).
\end{itemize}

\textbf{Reconocimiento:}
\begin{itemize}\setlength{\itemsep}{1pt}
  \item ``k-ésimo menor dinámico / posición de un elemento'' $\to$ \textit{ordered\_set}.
  \item ``\textit{unordered\_map} con claves numéricas (hackeadas)'' $\to$ \textit{custom\_hash}.
  \item ``$N \approx 36\dots42$: demasiado para $2^N$, sin subestructura DP'' $\to$ \textit{MeetInTheMiddle} (dos mitades de $2^{N/2}$, buscar con \textit{lower\_bound}).
  \item ``Comparar multiconjuntos de rangos en $O(1)$'' $\to$ \textit{ZobristHash} (prefijos XOR con hashes aleatorios por valor).
  \item ``Alcanzabilidad todos-pares en grafos densos ($N^3/64$)'' $\to$ \textit{BitsetReachability}.
  \item ``Comprimir coordenadas para segment tree/BIT/HLD'' $\to$ \textit{compress\_coords}.
\end{itemize}

\textbf{Errores y tips:}
\begin{itemize}\setlength{\itemsep}{1pt}
  \item \textit{ordered\_multiset}: inserta \textit{\{valor, id\}} para permitir duplicados; \textit{find\_by\_order} devuelve \textit{\{valor, id\}}.
  \item \textit{BitsetReachability} es 1-based.
  \item \textit{compress\_coords} devuelve ids 0-based en orden creciente de valor.
  \item \textit{MeetInTheMiddle.subset\_sum} asume valores no negativos para el objetivo exacto (verifica el problema).
\end{itemize}

---

\section*{11. Recomendaciones}
\begin{itemize}\setlength{\itemsep}{1pt}
  \item Tras cada competencia, actualiza la guía y la plantilla: si un problema requirió modificar un algoritmo, documenta la variación.
  \item Practica a conciencia los algoritmos que más se modifican en vivo: \textit{Dijkstra} (multi-estado), \textit{HLD}, segment trees (Node/merge y lazy), DP (identificación de recurrencia) y flujo (reducciones a min-cut).
  \item El inventario refleja el orden de \textit{plantilla.pdf}; si cambias los headers, recompila la plantilla y revisa que la guía siga alineada.
  \item Usa la guía como apoyo teórico y de navegación; el código vive en la plantilla.
\end{itemize}

\end{multicols}

\end{document}